\section{Discussion}
\label{sec:discussion}

\subsection{Relation to existing methods}

The framework places the major MIL methods in a common
assumption-naming language: each method corresponds to a particular
parametrization of, or surrogate for, the same coarsening problem.
\begin{itemize}
\item \textbf{Diverse Density} \citep{maron1998framework} is a
  noisy-OR maximum-likelihood specialization, and inherits the
  firing-rate confound of \cref{rem:firing-rate}.
\item \textbf{mi-SVM and MI-SVM} \citep{andrews2003support} are
  hard-assignment surrogates for the candidate-set sum under the C1
  (deterministic-OR) constraint.
\item \textbf{Attention-based deep MIL} \citep{ilse2018attention}
  learns a pooling map whose weights estimate instance scores only
  in the column space of the composition matrix
  (\cref{thm:rank,thm:bag-total}).
\item \textbf{CLAM} \citep{lu2021data} adds instance-level
  clustering supervision, which the framework reads as singleton
  candidate sets that improve the rank of the design.
\end{itemize}
These methods perform differently across applications
\citep{carbonneau2018multiple} for reasons the framework now
names: each method's pooling rule or surrogate implicitly assumes a
particular aggregation rule and a particular candidate-set
structure, and performs well when the assumption matches the data.

\subsection{Relation to partial-label learning}

Partial-label learning (PLL) \citep{cour2011learning} is the closest
identifiability-flavored neighbour of MIL and is worth a separate
comparison because its mathematical machinery is, in our framework,
a transposed instance of the same coarsening problem. In PLL, each
training instance carries a \emph{candidate label set} of size at
most $L$ that is known to contain the true label; the goal is to
learn an instance-level $L$-way classifier despite the per-instance
label ambiguity. In MIL, by contrast, each training bag carries a
\emph{candidate instance set} (the bag itself) together with one
observed binary label that is determined by the OR of latent
instance labels. The two settings exchange the roles of ``instance''
and ``label'' as the coarsened axis; both fit the masked-data
template of \cref{sec:background}.

The C1--C2--C3 conditions translate cleanly. C1 in PLL is the
guarantee that the true label is in the candidate set, the
defining assumption of the setting. C2 in PLL is the
condition that candidate-set composition does not depend on which
label is true given the instance features; \citet{cour2011learning}
formalize a sufficient quantitative version of this as the
\emph{small-ambiguity-degree} condition, under which the PLL
empirical risk minimizer is consistent. C3 is the assumption that
the candidate-set mechanism does not depend on the classifier
parameter. The rank-style identifiability machinery of
\cref{thm:rank} has a PLL analogue (a per-instance candidate-set
matrix replaces the bag composition matrix), and the bias bound of
\cref{thm:bias-rule} has a PLL analogue for misspecification of the
ambiguity mechanism. The two literatures have evolved
independently; the framework presented here makes their shared
identifiability content explicit, and suggests that imports from
the PLL side, particularly the small-ambiguity-degree analysis,
may sharpen finite-sample rates for MIL.

\subsection{What the framework does \emph{not} address}

\begin{itemize}
\item \textbf{Instance dependence within bags.} The DGP assumes
  instance labels are conditionally independent given type. Bags
  with correlated instances (non-i.i.d. MIL,
  \citealp{foulds2010review}) need a C2 relaxation.
\item \textbf{Fully nonparametric instance hazards.}
  \Cref{thm:rank-continuous} extends the identifiability machinery
  to continuous instance features under a linear-in-features
  parametrization $s(x) = \bm\beta^\top \bm\phi(x)$. A fully
  nonparametric treatment in which the per-instance hazard belongs
  to a more general function class would require RKHS or sieve
  techniques outside our scope.
\item \textbf{Deep-pooling identifiability.} Attention pooling may
  achieve practical identifiability through implicit regularization
  not captured by the linear rank condition. Whether such methods
  admit bias bounds comparable to \cref{thm:bias-rule} is open.
\item \textbf{Adaptive bag construction.} If bags are formed with
  knowledge of instance labels or of $\bm{\eta}$, C2 or C3 fails;
  the C2-relaxation results of \cite{towell2026mdrelax} would be
  the route to a quantitative treatment.
\end{itemize}

\subsection{Broader implications}

The framework shifts MIL analysis from method-specific to
assumption-specific reasoning. A practitioner asking ``can I trust
this model's instance-level outputs?'' can now answer by checking
whether the rank condition holds for the bag-composition structure,
whether the aggregation rule is plausibly deterministic OR, and
whether any instance-level labels are available to break the
firing-rate confound. The most actionable consequence is a caution:
by \cref{thm:bag-total}, a MIL model that predicts bag labels well
tells us nothing on its own about instance-level correctness, so
attention weights and instance scores require independent
instance-resolved validation. This is the same shift previously
achieved for scRNA-seq zero inflation
\citep{towell2026scrnacoarsening} and spatial-transcriptomics
deconvolution \citep{towell2026spatialcoarsening}.
